\honest{Feel the Meaning}{Eliezer Yudkowsky}{2008}

A friend says ``Oh, look, a butterfly.'' Pressure waves cross the air, my ear and nerves and cortex do their work, and within a second I know what my friend is looking at. If we did not know about the pressure waves, it would be front-page news: humans are telepathic. ``Well, we are telepathic, in fact.'' If you think it is simple, try building a computer that does it.

It would be inconvenient to think ``Now I shall partially transduce some features of my thoughts into a linear sequence of phonemes,'' so the brain hides this complexity, or never represents it. When a tiger leaps at me, I think ``Yikes! A tiger!'' and do not reason from its stripes. When someone shouts ``tiger,'' natural selection would not favor a listener who stopped to consider that the syllables ``Tie'' and ``Grr'' are what tribe members say when they see something they classify as, aiiieeee, CRUNCH. So there should be no step between hearing a word and activating its concept. In my blegg network I draw this as a direct link from the word ``Blegg!'' to the central unit, and for speaking, from the unit to shouting ``Blegg!'' \nb{This is an argument about what a good design would do, like the network in ``How an Algorithm Feels From Inside.'' It is not evidence about how brains are built.}

What that feels like from the inside is that the word and the concept are ``very nearly identified,'' and the meaning feels like a property of the word itself. \nb{C.~K. Ogden and I.~A. Richards wrote in 1923 that words ``mean'' nothing by themselves, ``as every one now knows.''} The cognoscenti will recognize E.~T. Jaynes's Mind Projection Fallacy: the meaning seems to belong to the word, as redness seems to belong to an apple or mysteriousness to a mysterious phenomenon. \nb{The term is explained in the post ``Mind Projection Fallacy,'' later in the book.}

On most occasions the brain does not separate word and meaning at all, ``only bothering to separate the two while learning a new language, perhaps,'' and even then, seeing Susan point and say ``Blegg!'', you wonder what ``blegg'' means, not what category Susan attaches to the sound. \nb{A familiar experience separates them: the word on the tip of the tongue, where a person has the meaning and cannot find the word. Brown and McNeill produced this state in 1966 by reading out definitions and asking for the words.}

Now back to Albert and Barry. Albert feels that ``sound'' has a meaning and that the meaning is vibrations, just as he feels the tree makes a sound, rather than causing an event that matches his category. Barry feels that its meaning is auditory experience, so the forest has no sound. In programmer's terms, Barry feels that sound.meaning is auditory experience and forest.sound is false. What is closer to the truth is that Barry's brain finds the concept it attaches to ``sound,'' and the forest does not match it. Humans did not evolve to know this, any more than they instinctively know the brain is made of neurons. Arguing over what a word means feels like arguing over a fact, like whether the sky is blue or green. You may not notice that anything is wrong until you try to state an experiment whose result depends on the dispute.
